% BHU PYQ -- Professional Exam Template
\documentclass[11pt,a4paper]{article}

\usepackage[a4paper,top=2.5cm,bottom=2.5cm,left=2.8cm,right=2.8cm,headheight=14pt]{geometry}
\usepackage{amsmath,amssymb}
\usepackage[T1]{fontenc}
\usepackage[utf8]{inputenc}
\usepackage{lmodern}
\usepackage{microtype}
\usepackage{fancyhdr}
\usepackage{enumitem}
\usepackage{hyperref}
\usepackage{lastpage}
\usepackage{parskip}

\hypersetup{colorlinks=true,urlcolor=black,linkcolor=black,
  pdftitle={B.Sc. (Hons.) Zoology Sem-I ZOB-101 -- BHU PYQ}}

\pagestyle{fancy}
\fancyhf{}
\renewcommand{\headrulewidth}{0.4pt}
\renewcommand{\footrulewidth}{0.4pt}
\fancyhead[L]{\small   \textit{Banaras Hindu University}}
\fancyhead[R]{\small   \textit{Zoology Paper: ZOB-101}}
\fancyfoot[C]{\small Page  \thepage\ of \pageref{LastPage}}


\newlist{parts}{enumerate}{2}
\setlist[parts,1]{label=(\alph*),leftmargin=*,itemsep=4pt,topsep=3pt}
\setlist[parts,2]{label=(\roman*),leftmargin=*,itemsep=2pt,topsep=2pt}

\newcommand{\pts}[1]{\hfill   \textit{[#1]}}


\begin{document}


\begin{center}
  {\large\bfseries BANARAS HINDU UNIVERSITY}\\[2pt]
  {
\normalsize B.Sc. (Hons.) Semester I Examination, 2023-24}\\[6pt]
  \rule{0.8\linewidth}{0.5pt}\\[6pt]
  {\large\bfseries Zoology}\\[4pt]
  {\large\bfseries Paper: ZOB-101 --- Animal Diversity and Animal Form and Function}\\[4pt]
  {
\normalsize Time: Three hours \hfill Full Marks: 70}
\end{center}

\rule{\linewidth}{0.4pt}

\smallskip



\noindent   \textit{Note: Answer three questions from Section-A (including Question No. 1 which is compulsory) and three questions from Section-B (including Question No. 1 which is compulsory). Use a separate answer book for each section. The figures in the right-hand margin indicate marks.}

\bigskip

\begin{center}
   \textbf{SECTION -- A} \\
   \textbf{(Animal Diversity) (Marks: 35)}
\end{center}




\noindent   \textbf{Question 1.} Answer the following parts:

\begin{parts}
  \item State whether the following statements are True or False. Give a brief justification for each:\pts{$2  \times 2 = 4$}
  
\begin{parts}
    \item Cnidoblasts are the stinging cells found in Ctenophora.
    \item   \textit{Archaeopteryx} is the connecting link between reptiles and birds.
  \end{parts}
  
  \item Define the following:\pts{$3  \times 1 = 3$}
  
\begin{parts}
    \item Amebocytes
    \item Phasmids
    \item Paedogenesis.
  \end{parts}
  
  \item Differentiate between the following:\pts{$2  \times 2 = 4$}
  
\begin{parts}
    \item Univalvia and Bivalvia
    \item Tentaculata and Nuda.
  \end{parts}
\end{parts}

\medskip


\noindent   \textbf{Question 2.} Answer the following:

\begin{parts}
  \item Explain why birds are called as glorified reptiles.\pts{6}
  \item Give a comparative account of adult and larval form of an ascidian (   \textit{Herdmania}).\pts{6}
\end{parts}

\medskip


\noindent   \textbf{Question 3.} Answer the following:

\begin{parts}
  \item Classify class Mammalia   \textit{or} Amphibia up to their orders giving salient features and examples of each.\pts{6}
  \item Classify phylum Mollusca   \textit{or} Echinodermata up to their classes giving characteristic features and suitable examples of each class.\pts{6}
\end{parts}

\pagebreak

\medskip


\noindent   \textbf{Question 4.} Write short notes on the following:\pts{$4  \times 3 = 12$}

\begin{parts}
  \item Parasitic adaptations in Platyhelminthes
  \item Neornithes
  \item Spicules in Porifera
  \item Retrogressive metamorphosis.
\end{parts}

\bigskip

\begin{center}
   \textbf{SECTION -- B} \\
   \textbf{(Animal Form and Function) (Marks: 35)}
\end{center}




\noindent   \textbf{Question 1.} Answer the following parts:

\begin{parts}
  \item State whether the following statements are True or False. Justify your answer in 2--3 sentences each:\pts{$4  \times 2 = 8$}
  
\begin{parts}
    \item Discontinuous feeders are active and mobile in nature.
    \item Freshwater fishes have aglomerular kidney.
    \item Complete separation of ventricle facilitates double circulation.
    \item The myelin sheath is an extension of a cell organelle.
  \end{parts}
  
  \item Define the following:\pts{$3  \times 1 = 3$}
  
\begin{parts}
    \item Deposit feeding
    \item Sequential hermaphroditism
    \item Venous heart.
  \end{parts}
\end{parts}

\medskip


\noindent   \textbf{Question 2.} Answer the following:

\begin{parts}
  \item With the help of diagrams, explain various methods of asexual reproduction. Comment on its advantages and disadvantages.\pts{6}
  \item Describe the feeding mechanisms in herbivorous and carnivorous animals citing suitable examples.\pts{6}
\end{parts}

\medskip


\noindent   \textbf{Question 3.} Answer the following:

\begin{parts}
  \item Describe the structure and function of fish gills. With the help of diagram, explain the mechanism of counter current gas exchange.\pts{6}
  \item Describe the structure and function of protonephridial excretory system with suitable diagrams and examples.\pts{6}
\end{parts}

\medskip


\noindent   \textbf{Question 4.} Write short notes on the following:\pts{$3  \times 4 = 12$}

\begin{parts}
  \item Otolith
  \item Cranial nerves
  \item Cephalization.
\end{parts}

\vfill
\rule{\linewidth}{0.4pt}

\begin{center}\small
\href{https://bhu-exam.vercel.app/}{Click Here}\\[4pt]\href{https://me-aryan.vercel.app/}{Click Here}\\[4pt]\href{https://quashan-me.vercel.app/}{Click Here}
\end{center}

\end{document}
